\documentclass[dvipsnames]{article}

\usepackage[utf8]{inputenc}
\usepackage{fontspec}
\usepackage{polyglossia}
\usepackage{unicode-math}
\usepackage{csquotes}
\setdefaultlanguage{russian}
\setotherlanguage{english}

\setmainfont{CMU Serif}
\setsansfont{CMU Sans Serif}
\setmonofont{CMU Typewriter Text}
% \setmathfont{Latin Modern Math}

\usepackage[letterpaper,top=2cm,bottom=2cm,left=3cm,right=3cm,marginparwidth=1.75cm]{geometry}

\usepackage{amsmath}
\usepackage{graphicx}
\usepackage{amsthm}
\usepackage{array}
\usepackage{fancyhdr}
\usepackage{bussproofs}
\usepackage[colorlinks=true, allcolors=blue]{hyperref}
\usepackage[backend=biber]{biblatex}
\usepackage{xcolor}

\usepackage{minted2}
\usemintedstyle{xcode}
\setminted{escapeinside=||}
\newminted[stella]{cpp}{escapeinside=||,frame=single}
\newmintinline[haskell]{haskell}{fontsize=\scriptsize}
\newmintinline[java]{java}{fontsize=\scriptsize}
\newmintinline[python]{python}{fontsize=\scriptsize}
\newmintinline[code]{c}{fontsize=\small}
\newmintinline[smallcode]{c}{fontsize=\small}
\let\smallcodeOrig\smallcode
\renewcommand{\smallcode}[2][]{%
  \ifmmode
    \if\relax\detokenize{#1}\relax
      \mathchoice{\mbox{\normalfont\small\ttfamily #2}}{\mbox{\normalfont\small\ttfamily #2}}{\mbox{\normalfont\small\ttfamily #2}}{\mbox{\normalfont\small\ttfamily #2}}%
    \else
      \mathchoice{\mbox{\normalfont\small\ttfamily #2}}{\mbox{\normalfont\small\ttfamily #2}}{\mbox{\normalfont\small\ttfamily #2}}{\mbox{\normalfont\small\ttfamily #2}}%
    \fi
  \else
    \smallcodeOrig[#1]{#2}%
  \fi}
\let\codeOrig\code
\renewcommand{\code}[2][]{%
  \ifmmode
    \if\relax\detokenize{#1}\relax
      \mathchoice{\mbox{\normalfont\small\ttfamily #2}}{\mbox{\normalfont\small\ttfamily #2}}{\mbox{\normalfont\small\ttfamily #2}}{\mbox{\normalfont\small\ttfamily #2}}%
    \else
      \mathchoice{\mbox{\normalfont\small\ttfamily #2}}{\mbox{\normalfont\small\ttfamily #2}}{\mbox{\normalfont\small\ttfamily #2}}{\mbox{\normalfont\small\ttfamily #2}}%
    \fi
  \else
    \codeOrig[#1]{#2}%
  \fi}

\usepackage{sourcecodepro}
\usepackage{mdframed}

\DeclareMathOperator{\xor}{\mathsf{xor}}
\DeclareMathOperator{\tru}{\mathsf{tru}}
\DeclareMathOperator{\fls}{\mathsf{fls}}

\addbibresource{../references.bib}

\begin{document}

\pagestyle{fancy}
\fancyhf{}
\chead{Задание выполнил \textcolor{blue}{Павел Григоренко}}
\cfoot{Языки программирования. Семантика и системы типов (весна 2026, Университет ИТМО)}

\section*{Тема 5. Теоретическое задание \textcolor{blue}{(выполнил Павел Григоренко)}}%

Расширьте доказательство теоремы о нормализации простого типизированного $\lambda$-исчисления,
добавив поддержку \code{let}-связываний~\cite[Рис.~11.4]{TaPL-ru}, типов-записей~\cite[Рис.~11.7]{TaPL-ru} и типов-сумм~\cite[Рис.~11.9]{TaPL-ru}:

\begin{enumerate}
  \item Расширьте определение $R_{\smallcode{T}}(\code{t})$ для случаев
  \begin{itemize}
    \item $\code{T} = \{\code{l$_1$ : T$_1$, …, l$_n$ : T$_n$}\}$
    \item $\code{T} = \code{T}_1\code{ + T}_2$
  \end{itemize}

  \color{blue}
  \textbf{Решение.}
  \emph{место для решения}
  \color{black}

  \clearpage

  \item Расширьте доказательство леммы о сохранении свойства $R_{\smallcode{T}}$~\cite[Лемма~12.1.4]{TaPL-ru}.

  \color{blue}
  \textbf{Решение.}
  \emph{место для решения}
  \color{black}

  \clearpage

  \item Расширьте доказательство леммы о нормализации открытых термов~\cite[Лемма~12.1.5]{TaPL-ru},
  рассмотрев следующие правила типизации:
  \begin{itemize}
    \item \textsc{T-Rcd}, \textsc{T-Proj},
    \item \textsc{T-Inl}, \textsc{T-Inr}, \textsc{T-Case}
    \item \textsc{T-Let}
  \end{itemize}

  \color{blue}
  \textbf{Решение.}
  \emph{место для решения}
  \color{black}
\end{enumerate}

\section*{Список литературы}
\printbibliography[heading=none]

\end{document}

